% ================= 5. PRELIMINARY HARDWARE DESIGN =================
\section{Preliminary Hardware Design}
\label{sec:hardware}

This design is preliminary: final actuator and sensor specifications are not yet
confirmed, so the mechanical and electrical architecture is kept modular and
largely parametric (Section~\ref{sec:background}).

\subsection{System Architecture}
\label{sec:architecture}

% High-level block diagram tying together actuation, sensing, processing, and
% power for both the 2D prototype and the 3D cube.

\subsection{Mass and Inertia Budget}
\label{sec:massbudget}

Before committing to component geometry, the preliminary design is bounded by
three coupled constraints that must be satisfied simultaneously. Sizing to the
worst case of each fixes the envelope within which the mechanical and electrical
design then iterate.

\paragraph{Volume constraint.} Every subsystem---the three orthogonal
reaction-wheel/motor modules, the battery, and the electronics stack---must fit
inside the $180$~mm cube envelope. The binding dimension is the reaction wheel:
its outer diameter must clear the internal frame and corner/edge contact
features while leaving room for the motor and the servo brake. This caps the
usable wheel diameter and, with it, the radius at which ballast mass can be
placed---directly limiting the inertia achievable at a given wheel mass
(Section~\ref{sec:rw_sizing}).

\paragraph{Mass constraint.} The cube is designed to a total mass target of
$M = 1.45$~kg. This is not merely a logistics figure: it feeds back into the
dynamics, since the gravitational tip-over torque $\tau_g = MgL/2$ and the
required jump-up angular momentum both scale with $M$ (Section~\ref{sec:jumpup_target}).
A heavier cube demands a more capable brake and a larger wheel inertia, so the
mass target and the inertia target are coupled and must be closed together. The
preliminary allocation is given in Table~\ref{tab:massbudget}, with the
reaction-wheel set (the one firmly-sized item) committed and the remaining
subsystems held as allowances to be refined as components are selected.

\begin{table}[h]
\centering
\caption{Preliminary mass budget (per cube). Reaction wheels sized; others allocated.}
\label{tab:massbudget}
\small
\begin{tabular}{l S[table-format=4.0] l}
\toprule
Subsystem & {Allocated mass (g)} & Status \\
\midrule
Reaction wheels ($3\times$)   & 395   & Sized (Sec.~\ref{sec:rw_sizing}) \\
Motors ($3\times$)            & {TBD} & Allowance --- pending motor selection \\
Frame / structure             & {TBD} & Allowance --- pending CAD \\
Battery                       & {TBD} & Allowance --- pending capacity sizing \\
Electronics (MCU, drivers, IMU, servo) & {TBD} & Allowance \\
Wiring / fasteners / margin   & {TBD} & Allowance \\
\midrule
\textbf{Total (target)}       & \bfseries 1450 & \textbf{Design target} \\
\bottomrule
\end{tabular}
\end{table}

\paragraph{Inertia constraint (corner-critical).} Each reaction wheel must store
enough angular momentum at the motor's maximum speed to execute the jump-up
maneuver. The two contact cases are not equally demanding: the corner
equilibrium requires a larger rise of the center of mass and a less favorable
momentum-transfer geometry than the edge equilibrium, so for a common maximum
wheel speed $\omega_{\max}$ it demands the greater wheel inertia. Evaluating the
torque-limited budget of Section~\ref{sec:jumpup_target} for both cases
(Table~\ref{tab:inertia_cases}), the corner case requires
$3.15\times10^{-4}$ kg\,m$^2$ against $2.85\times10^{-4}$ kg\,m$^2$ for the edge
---about $10\%$ more. The corner jump-up is therefore the critical condition,
and the reaction wheel is sized to it; a wheel that satisfies the corner
requirement satisfies the edge requirement with margin.

\begin{table}[h]
\centering
\caption{Required per-wheel inertia by contact case, at $\omega_{\max}=8400$~rpm.}
\label{tab:inertia_cases}
\begin{tabular}{l S[table-format=1.3] S[table-format=1.2]}
\toprule
Contact case & {$h_w$ (kg\,m$^2$/s)} & {$I_{w,\text{target}}$ ($10^{-4}$ kg\,m$^2$)} \\
\midrule
Edge                       & 0.251 & 2.85 \\
\textbf{Corner (critical)} & \bfseries 0.277 & \bfseries 3.15 \\
\bottomrule
\end{tabular}
\end{table}

\subsection{Mechanical Design}
\label{sec:mechanical}

\subsubsection{2D Prototype Build}

% Structural concept of the single-face rig: 3D-printed plate, central
% motor/wheel mount, corner bearing, and servo brake (cross-reference
% Sec.~\ref{sec:2dmodel}).

\subsubsection{3D Cube Frame}

% Structural concept for the full cube: three orthogonal wheel/motor modules,
% frame, and corner/edge contact features.

\subsubsection{Reaction-Wheel Sizing}
\label{sec:rw_sizing}

The reaction wheel must supply the per-wheel inertia
$I_{w,\text{target}} = 3.15\times10^{-4}$ kg\,m$^2$ derived from the jump-up
budget of Section~\ref{sec:jumpup_target}, while remaining light enough not to
dominate the cube mass. Rather than adjusting inertia by scaling a solid disc,
the wheel is built as a printed \emph{ring plus spokes} whose inertia is then
tuned by \emph{ballast}: standard steel hex nuts dropped into hexagonal pockets
distributed around the ring. Each pocket both removes plastic and adds steel, so
the net inertia gain per pocket is large and the design target is reached by
choosing how many nuts to install---a discrete, easily-adjustable knob that also
lets the physical wheel be re-tuned after fabrication to absorb parameter
uncertainty.

The structure is printed in PET-CF ($\rho = 1290$ kg/m$^3$): a
$120$~mm outer-diameter ring of $15$~mm radial width and $10$~mm thickness,
carried by three radial spokes $10$~mm wide of the same thickness. The hub and
motor interface are neglected in the inertia budget. The ballast is standard
M5 steel hex nuts ($\rho = 7850$ kg/m$^3$, $11$~mm across flats,
measured mass $3.12$~g), seated in blind hexagonal pockets machined at the
mid-width of the ring. Accounting for the plastic each pocket displaces, the
\emph{net} contribution per installed nut is $+2.39$~g of mass and
$+1.55\times10^{-5}$ kg\,m$^2$ of inertia at the ring radius.

Table~\ref{tab:rw_sizing} sweeps the outer diameter, solving for the number of
nuts needed to reach $I_{w,\text{target}}$ and checking that the pockets
physically fit. Small wheels need so many nuts that no plastic wall survives
between adjacent pockets; wheels of $140$~mm and above reach the target from the
bare ring alone (no tuning margin, and heavier). The selected design is
$\mathbf{OD = 120}$~\textbf{mm with 21 nuts}: it meets the inertia target with a
$\sim\!2.5$~mm circumferential wall between pockets, splits the inertia roughly
evenly between structure ($45\%$) and ballast ($55\%$) so the wheel can be
trimmed in either direction, and keeps the three-wheel set to $\sim\!395$~g
($27\%$ of the $1.45$~kg cube).

\begin{table}[h]
\centering
\caption{Reaction-wheel outer-diameter trade study. Selected design in bold.}
\label{tab:rw_sizing}
\small
\setlength{\tabcolsep}{4pt}
\begin{tabular}{S[table-format=3.0] S[table-format=3.0] S[table-format=3.1] S[table-format=1.2] S[table-format=3.1] S[table-format=3.1] l}
\toprule
{OD (mm)} & {$N$} & {Mass (g)} & {$I_w$ ($10^{-4}$\,kg\,m$^2$)} & {$I_w$/target (\%)} & {3 wheels (g)} & {Pocket fit} \\
\midrule
80  & 105 & 300.6 & 3.17 & 100.7 & 901.7 & fails --- no wall \\
90  & 72  & 229.6 & 3.16 & 100.5 & 688.7 & fails --- no wall \\
100 & 51  & 187.3 & 3.25 & 103.3 & 561.9 & fails --- no wall \\
110 & 33  & 152.2 & 3.22 & 102.3 & 456.7 & fails --- thin wall \\
\bfseries 120 & \bfseries 21 & \bfseries 131.5 & \bfseries 3.31 & \bfseries 105.2 & \bfseries 394.6 & \textbf{OK ($\sim$2.5\,mm)} \\
130 & 9   & 110.8 & 3.23 & 102.7 & 332.4 & OK (loose) \\
140 & 0   & {---} & 3.23 & {---} & {---} & bare ring exceeds target \\
150 & 0   & {---} & 4.07 & {---} & {---} & bare ring exceeds target \\
\bottomrule
\end{tabular}
\end{table}

\paragraph{Radial fit caveat.} The circumferential fit is satisfied at
$120$~mm, but the \emph{radial} fit is marginal: an $11$~mm across-flats pocket
in a $15$~mm-wide ring leaves only $1.8$~mm of wall on each side (dropping to
$0.9$~mm if a nut is seated corner-out), below the $2$~mm minimum-wall target.
The ring width should therefore be increased to $\gtrsim\!17$~mm, or the nut
orientation constrained flats-radial, in the next design iteration. This does
not change the inertia budget, only the ring's structural robustness around the
pockets.

\subsection{Electrical Design}
\label{sec:electrical}

\subsubsection{Components}

The electronics and software will be built around a microcontroller communicating over a CAN bus with a motor driver that closes the current loop for the selected brushless motor. Wheel speed will be measured with a magnetic encoder on the motor shaft, and body tilt will be estimated from an onboard IMU. Braking will be actuated by a digital servo driven directly from the microcontroller, and a wireless module will provide a telemetry and debugging link. Power will be supplied by a LiPo battery, stepped down through a DC-DC converter for the logic-level electronics, with the motor driver powered directly from the battery bus.

The full parts list is catalogued in Appendix~\ref{app:bom}
(Table~\ref{tab:bom-supplied}). The supplied set is an electronics-only BOM;
the reaction wheels, the frame and all printed parts are the team's design
responsibility, and their fabrication is the leading schedule risk.

\subsubsection{Electrical Diagram}

% Wiring / power-distribution diagram: battery bus, DC-DC converter, motor
% driver, microcontroller, encoder, IMU, servo, and telemetry link.

\subsection{System Integration}
\label{sec:integration}

% How the actuation, sensing, processing, and power subsystems are assembled and
% coordinated: mechanical mounting, interface management, and the software
% coordination that runs the control loop derived in Sec.~\ref{sec:modeling}.
